\section{Điều gradient chưa giải quyết}
\label{sec:ch05-dieu-gradient-chua-giai-quyet}

Grad-CAM cho một bản đồ có điều kiện theo lớp, nhưng bản đồ ấy không phải chứng
minh rằng mô hình dùng đúng vùng mà con người mong đợi. Nó là một phép tổng hợp
từ gradient và biểu diễn của một lớp được chọn. Nếu lưới ở lớp ấy thô, vùng
được tô cũng thô. Nếu mô hình dựa vào nền ảnh hoặc một tương quan không mong
muốn, Grad-CAM có thể làm lộ điều đó, nhưng cũng không tự phân biệt tương quan
với quan hệ nhân quả.

Lý do tương tự áp cho Integrated Gradients. Một phương pháp có thể thỏa \tn{độ nhạy}{sensitivity},
\tn{bất biến theo cài đặt}{implementation invariance} và \tn{tính đầy đủ}{completeness}, song vẫn trả lời một phản thực vô nghĩa
nếu baseline không phù hợp. Một lỗi thường gặp là biến các bảo đảm hình thức
thành kết luận về \tn{độ trung thực}{faithfulness}. Các điều kiện hình thức
nói attribution nhất quán với hàm trên đường đi đã chọn; kết luận về độ trung
thực đòi bằng chứng rằng đường đi, điểm nền hoặc lớp biểu diễn tương ứng với
thay đổi có nghĩa trong miền. Với
Integrated Gradients, \tn{điểm nền}{baseline} là phần của lựa chọn ấy ngay từ
lúc xác định câu hỏi.

Hai bài báo vì thế để lại cùng một công việc cho các chương sau. Chương
\ref{ch:attribution-tu-nguyen-ly-dau} sẽ hỏi attribution nên được xây từ những
nguyên lý nào trước khi chọn một phương pháp. Từ chương~\ref{ch:do-do-trung-thuc},
cuốn sách quay sang các chỉ số dùng để đánh giá những lời giải thích ấy. Khi
đó, ta sẽ không thể chỉ kiểm tra một bản đồ trông hợp lý; chính dụng cụ đo độ
trung thực cũng phải được kiểm tra.

\index{Grad-CAM|)}%
\index{độ trung thực}%
